\section{PECL Specialization: Support Reduction onto $\mathcal{W}_{\Delta,t}$}
\label{sec:pecl-spec}

Theorem~\ref{thm:pathwise_pac} holds for any hypothesis class and any distribution family for the posterior.
We now specialize it to frozen-backbone PECL by showing that both the KL complexity and the sharpness-dependent empirical term reduce exactly to quantities supported on the task-admissible subspace $\mathcal{W}_{\Delta,t}$.

\subsection{KL Reduction}
\label{sec:kl-reduction}

In PECL, the task-$t$ parameters decompose as $W_t = W_t^{\mathrm{froz}} + \Delta W_t$ with $\Delta W_t \in \mathcal{W}_{\Delta,t}$.
We place a Gaussian posterior over increments:
\[
Q_t^\Delta = \mathcal{N}(\widehat{\Delta W}_t, \Sigma_t),\quad \Sigma_t \text{ supported on } \mathcal{W}_{\Delta,t}.
\]
Sampling $W_t \sim Q_t$ on the affine class $W_t^{\mathrm{froz}} + \mathcal{W}_{\Delta,t}$ amounts to setting $W_t = W_t^{\mathrm{froz}} + \Delta W_t$ with $\Delta W_t = \widehat{\Delta W}_t + \varepsilon$ and $\varepsilon \sim \mathcal{N}(0, \Sigma_t)$.

\begin{figure}[t]
\centering
\includegraphics[width=0.9\linewidth]{figures_TRML/KLshiyi1.pdf}
\caption{\textbf{KL reduction.}
The posterior/prior defined on task-permissible subspaces $(Q_t^\Delta, P_t^\Delta)$ induce full-model distributions $(Q_t, P_t)$ satisfying $\mathrm{KL}(Q_t \| P_t) = \mathrm{KL}(Q_t^\Delta \| P_t^\Delta)$.
The frozen backbone contributes only degenerate factors that cancel in the KL ratio, so full-model complexity depends only on the admissible subspace.}
\label{fig:KL_decompose}
\end{figure}

\begin{lemma}[KL reduction under frozen backbone]
\label{lem:kl-decomp-pecl}
Assume $Q_t^\Delta \ll P_t^\Delta$. Then $Q_t \ll P_t$ and
\[
\mathrm{KL}(Q_t \| P_t) = \mathrm{KL}(Q_t^\Delta \| P_t^\Delta).
\]
\end{lemma}

This follows because the frozen backbone contributes identical (degenerate) factors to both $Q_t$ and $P_t$, which cancel in the Radon–Nikodym derivative.
The full-model KL divergence is therefore determined entirely by the distribution of the admissible increment $\Delta W_t$, not by any frozen directions.
See Appendix~\ref{proof:Lemma1} for the proof.

\subsection{Subspace Sharpness Decomposition}
\label{sec:sharpness-decomp}

We define the \emph{subspace sharpness} as the expected loss elevation under perturbations supported on $\mathcal{W}_{\Delta,t}$:
\[
\mathrm{TS}_t^\Delta(\widehat{\Delta W}_t, \Sigma_t) = \mathbb{E}_{\varepsilon \sim \mathcal{N}(0,\Sigma_t)} \bigl[\hat{L}_{S_t}(W_t^{\mathrm{froz}} + \widehat{\Delta W}_t + \varepsilon) - \hat{L}_{S_t}(W_t^{\mathrm{froz}} + \widehat{\Delta W}_t)\bigr],
\]
where $\varepsilon \in \mathcal{W}_{\Delta,t}$.

\begin{lemma}[Subspace sharpness decomposition]
\label{lem:subspace-loss-pecl}
For each task $t$,
\begin{align*}
\mathbb{E}_{W_t \sim Q_t}\bigl[\hat{L}_{S_t}(W_t)\bigr]
&= \hat{L}_{S_t}(W_t^{\mathrm{froz}} + \widehat{\Delta W}_t) + \mathrm{TS}_t^\Delta(\widehat{\Delta W}_t, \Sigma_t).
\end{align*}
\end{lemma}

The expected loss under the Gaussian posterior equals the deterministic loss at the learned increment plus the subspace sharpness correction.
Crucially, the sharpness term is supported entirely on $\mathcal{W}_{\Delta,t}$: the perturbation $\varepsilon$ probes only directions within the admissible subspace.
See Appendix~\ref{proof:lemma2} for the proof.

\subsection{PAC-Bayesian Bound for PECL}
\label{sec:pecl-bound}

Substituting Lemma~\ref{lem:kl-decomp-pecl} and Lemma~\ref{lem:subspace-loss-pecl} into Theorem~\ref{thm:pathwise_pac}, we obtain the main result:

\begin{theorem}[PAC-Bayesian bound for PECL]
\label{thm:pecl-bound}
Consider a PECL model with frozen backbone $W_t^{\mathrm{froz}}$ and Gaussian posteriors $Q_t^\Delta = \mathcal{N}(\widehat{\Delta W}_t, \Sigma_t)$ supported on $\mathcal{W}_{\Delta,t}$.
Let $P_t^\Delta$ be task priors sampled from a hyperposterior $\mathcal{P}_{t-1}$ over adapter distributions.
For any $\delta \in (0,1]$, with probability at least $1 - \delta$ over $S_{1:T}$,
\begin{equation*}
\resizebox{\linewidth}{!}{$
\begin{aligned}
&\frac{1}{T}\sum_{t=1}^T \mathbb{E}_{P_t^\Delta \sim \mathcal{P}_{t-1}} \mathbb{E}_{W_t \sim Q_t} L_{\mathcal{D}_t}(W_t) \\
&\leq
\frac{1}{T}\sum_{t=1}^T \mathbb{E}_{P_t^\Delta \sim \mathcal{P}_{t-1}}
\bigl[
\hat{L}_{S_t}(W_t^{\mathrm{froz}} + \widehat{\Delta W}_t)
+ \mathrm{TS}_t^\Delta(\widehat{\Delta W}_t, \Sigma_t)
\bigr] \\
&\quad +
\sqrt{
\frac{
\displaystyle\sum_{t=1}^T \mathbb{E}_{P_t^\Delta \sim \mathcal{P}_{t-1}}
\mathrm{KL}(Q_t^\Delta \| P_t^\Delta)
+ \sum_{t=1}^T \mathrm{KL}(\mathcal{P}_t \| \mathcal{P}_{t-1})
+ \log\tfrac{1}{\delta}
}{2\displaystyle\sum_{t=1}^T (m_t - 1)}
}.
\end{aligned}
$}
\end{equation*}
\end{theorem}

See Appendix~\ref{proof:theorem} for the proof.

\textbf{Interpretation: adapter-subspace flatness is the bound-relevant notion.}
Theorem~\ref{thm:pecl-bound} shows that, under frozen-backbone PECL, every quantity that governs the bound is supported on the admissible subspaces $\{\mathcal{W}_{\Delta,t}\}_{t=1}^T$.

The \emph{KL complexity} depends on $\mathrm{KL}(Q_t^\Delta \| P_t^\Delta)$, the divergence between prior and posterior on admissible adapter increments.
The frozen backbone contributes only a fixed translation and does not enter this term as an independently controllable complexity source.

The \emph{sharpness-smoothed empirical term} involves $\mathrm{TS}_t^\Delta$, the expected loss elevation under perturbations in $\mathcal{W}_{\Delta,t}$.
Curvature along frozen backbone directions is not directly governed by this term: those directions lie outside the admissible perturbation support.

This yields the precise claim: in frozen-backbone PECL, admissible-update flatness is the \emph{bound-relevant} notion of flatness.
This does not imply that full-space flatness is uninformative---it may correlate with adapter-subspace flatness, and empirical full-space probes may still reflect consequences of subspace-level flatness control.
But full-space flatness is not the quantity directly governed by the PECL bound.

\textbf{Perturbation scope follows from the bound.}
Since the bound depends only on perturbations within $\mathcal{W}_{\Delta,t}$, perturbing the frozen backbone introduces noise along directions that contribute nothing to the bound and cannot be compensated by subsequent LoRA updates.
This provides a theoretical rationale for restricting perturbation scope to the adapter subspace in PECL.

\textbf{Perturbation type within the subspace.}
Let $\Pi_{\Delta,t}$ be the projector onto $\mathcal{W}_{\Delta,t}$ and $C_{\Delta,t} = \Pi_{\Delta,t} C_t \Pi_{\Delta,t}$ be the restricted curvature operator.
Under a local second-order approximation,
\[
\mathrm{TS}_t^\Delta(\widehat{\Delta W}_t, \Sigma_t) \approx \tfrac{1}{2}\operatorname{tr}(C_{\Delta,t} \Sigma_t),
\]
where $\Sigma_t$ is the perturbation covariance supported on $\mathcal{W}_{\Delta,t}$.
Different perturbation types (RS$^{(0)}$, AS$^{(0)}$, AS$^{(1)}$) correspond to different choices of $\Sigma_t$: isotropic, adversarially concentrated, or gradient-aligned.
Type AS$^{(1)}$ aligns perturbation mass with the dominant curvature directions, minimizing $\operatorname{tr}(C_{\Delta,t}\Sigma_t)$ most directly and thus exerting the strongest control on the bound.
